%-------------------------------------------------------------------------------
%	SECTION TITLE
%-------------------------------------------------------------------------------
\cvsection{项目经历}


%-------------------------------------------------------------------------------
%	CONTENT
%-------------------------------------------------------------------------------
\begin{cventries}

%---------------------------------------------------------
  \cventry
    {容器化部署与运维}
    {AI 数字人直播系统（边缘单机部署）}
    {}
    {2025.06 - 2026.07}
    {
      \begin{cvitems}
        \item {\textbf{项目内容}：面向直播带货的 AI 数字人交互系统，包含本地 LLM 推理、语音合成与直播弹幕接入链路，需在单机边缘环境内稳定运行并接入商品与话术资料库。}
        \item {\textbf{项目职责}：独立完成 LLM 推理、GPT-SoVITS 语音链路与直播弹幕接入的部署；使用 Docker 与 docker-compose 管理服务依赖、启动顺序与重启策略；使用 Ollama 部署 Qwen3.5 27B（4.151 bpw 量化），挂载商品 / 话术资料库并完成 LoRA 促单话术适配；编写 Shell / Python 脚本处理服务启动、状态检查、日志采集与进程恢复；完成连续 6 小时运行验证，端到端交互约 17 秒。}
      \end{cvitems}
    }

%---------------------------------------------------------
  \cventry
    {Linux 系统构建与开源贡献}
    {多平台 OpenWrt / Armbian 系统移植与内核主线化}
    {}
    {2024.09 - 至今}
    {
      \begin{cvitems}
        \item {\textbf{项目内容}：基于 \mbox{OpenWrt} / \mbox{Armbian} 的多平台嵌入式 Linux 系统适配与维护，覆盖 TH1520、SG2000、JH7110、MT798x、Siflower 等 SoC 及 RISC-V64 / ARM64 / MIPS / x86 平台，包含系统配置、交叉编译、内核配置、Device\nobreakspace{}Tree 编译与镜像打包全流程。}
        \item {\textbf{项目职责}：跟踪上游内核驱动变更，使用 QEMU / GDB、串口输出与系统日志复现并定位启动、内存管理与驱动问题；参与 Siflower sf21 / sf19a2890 内核由 6.12 升级至 6.18，分析 DPNS 以太网驱动与 DMA / NPI 子系统相关修复；向 \mbox{OpenWrt} 主线合入 3 个内核补丁，为 QModem 增加高通 SIM8202G-M2 5G 模块支持，维护 \mbox{ImmortalWrt} user-FAQ 文档。}
      \end{cvitems}
    }

%---------------------------------------------------------
  \cventry
    {毕业设计}
    {RISC-V 平台智能家居 AI 助手}
    {}
    {2025.09 - 2026.06}
    {
      \begin{cvitems}
        \item {\textbf{项目内容}：在 RISC-V / \mbox{Armbian} 设备上构建的本地智能家居 AI 助手，整合语音交互、多轮对话与智能家居设备控制，全部组件运行于本地硬件。}
        \item {\textbf{项目职责}：在 RTX 3080 Laptop（16GB 显存）使用 Ollama 部署 DeepSeek-R1-Distill-Qwen-14B 作为对话后端；整合 SillyTavern、GPT-SoVITS 与 Home Assistant，实现语音交互、多轮对话与智能家居控制；使用 Docker 管理服务运行依赖，配置组件网络通信与启动脚本；使用 Coding Agent 辅助编码并逐项人工验证。}
      \end{cvitems}
    }

%---------------------------------------------------------
  \cventry
    {容器化建站与嵌入式工具}
    {LNMP 网站环境与嵌入式系统开发}
    {}
    {2022 - 2025}
    {
      \begin{cvitems}
        \item {\textbf{项目内容}：基于 Docker 的 LNMP 全栈网站环境与嵌入式工具链，覆盖公司及个人网站的建站上线与嵌入式固件开发。}
        \item {\textbf{项目职责}：搭建 LNMP（Nginx / MySQL / PHP）环境并完成性能调优与安全加固，处理域名解析、HTTPS 配置与站点日常维护；开发维护基于 STM32 的内存条 SPD 烧录器工具链及固件；完成基于树莓派 4 的竞速越野小车硬件电路设计与视觉识别系统部署。}
      \end{cvitems}
    }

%---------------------------------------------------------
\end{cventries}
